%%======================================================================
%%  MR-7 Literature Review:
%%  Graviton Scattering Amplitudes in Nonlocal and Higher-Derivative Gravity
%%======================================================================
\documentclass[11pt,a4paper]{article}

%% ---- packages -------------------------------------------------------
\usepackage[utf8]{inputenc}
\usepackage[T1]{fontenc}
\usepackage{lmodern}
\usepackage{amsmath,amssymb,amsthm,mathtools}
\usepackage[margin=2.5cm]{geometry}
\usepackage{booktabs}
\usepackage{enumitem}
\usepackage{hyperref}
\usepackage{xcolor}
\usepackage{fancyhdr}
\usepackage{longtable}

%% ---- theorem-like environments
\newtheorem{theorem}{Theorem}[section]
\theoremstyle{remark}
\newtheorem{remark}[theorem]{Remark}

%% ---- custom commands
\newcommand{\dd}{\mathrm{d}}
\newcommand{\ii}{\mathrm{i}}
\newcommand{\erfi}{\operatorname{erfi}}
\DeclareMathOperator{\PV}{P.V.}

%% ---- page style
\pagestyle{fancy}
\fancyhf{}
\fancyhead[L]{\small\textsc{MR-7 Literature: Graviton Scattering}}
\fancyhead[R]{\small\thepage}
\renewcommand{\headrulewidth}{0.4pt}

%% ---- hyperref
\hypersetup{
  colorlinks=true,
  linkcolor=blue!60!black,
  citecolor=green!50!black,
  urlcolor=blue!70!black,
  pdftitle={MR-7 Literature Review},
  pdfauthor={David Alfyorov},
}

%% ---- title
\title{\textbf{MR-7 Literature Review}\\[8pt]
  \large Graviton Scattering Amplitudes in Nonlocal\\
  and Higher-Derivative Gravity}

\author{David Alfyorov}
\date{March 2026}

\begin{document}

\maketitle

\tableofcontents

%%======================================================================
\section{Field Redefinition Theorem}
%%======================================================================

\subsection{Modesto--Calcagni (2021)}

The central result for MR-7 is the field redefinition theorem of
Modesto and Calcagni~\cite{ModestoCalcagni2021}
(arXiv:2107.04558, JHEP). For a nonlocal theory of the form
\begin{equation}
  S(\Phi) = S_{\mathrm{loc}}(\Phi) + E_i\, F_{ij}(\Delta)\, E_j + V(E),
\end{equation}
where $E_i = \delta S_{\mathrm{loc}}/\delta\Phi_i$, all on-shell
tree-level $n$-point amplitudes equal those of $S_{\mathrm{loc}}$.

Two independent proof methods are given:
\begin{enumerate}
  \item \textbf{Field redefinition:} $\Phi' = \Phi + Q_{ij} E_j$ maps
        $S(\Phi) = S_{\mathrm{loc}}(\Phi')$, preserving on-shell amplitudes.
  \item \textbf{Perturbative EoM:} Perturbative solutions of the nonlocal
        EoM satisfy the local EoM at every order.
\end{enumerate}

Application to SCT requires three conditions, all verified:
(1)~$E \cdot F \cdot E$ form with $V = 0$,
(2)~$F_1, F_2$ entire (NT-2),
(3)~background is a vacuum Einstein solution.

\subsection{Dona--Giaccari--Modesto--Rachwal--Zhu (2015)}

Explicit computation of the 4-graviton tree-level amplitude in
$d = 4, 5, 6$ for nonlocal gravity~\cite{DonaGiaccari2015}
(arXiv:1506.04589). Amplitudes are identical to GR when
$\gamma_4(\Box) = 0$ (no Riemann-squared term). In $d=4$, the
Gauss--Bonnet identity guarantees this condition is automatically
satisfied.

%%======================================================================
\section{One-Loop UV Behavior}
%%======================================================================

\subsection{Anselmi--Piva (2018): Fakeon One-Loop}

Anselmi and Piva~\cite{AnselmiPiva2018a} (arXiv:1803.07777)
computed the one-loop graviton self-energy (2-point function) in
Stelle gravity with the spin-2 ghost quantized as a fakeon. Key
results:
\begin{itemize}
  \item The absorptive part is controlled by the central charge $c$
        of the coupled matter.
  \item The fakeon prescription disentangles real parts (renormalizable
        power counting) from imaginary parts (low-energy expansion).
  \item The theory is one-loop renormalizable AND unitary.
\end{itemize}

The SCT analogue has entire-function (not polynomial) form factors
and infinitely many ghost poles, requiring extension of the Anselmi--Piva
framework.

\subsection{'t Hooft--Veltman (1974) and Goroff--Sagnotti (1986)}

't~Hooft and Veltman~\cite{tHooftVeltman1974} showed that pure gravity
is one-loop finite on-shell (off-shell divergences are proportional to the
equations of motion). Goroff and Sagnotti~\cite{GoroffSagnotti1986}
showed that two-loop gravity is divergent with the counterterm
$R_{\mu\nu\rho\sigma} R^{\rho\sigma\alpha\beta} R_{\alpha\beta}^{\ \ \mu\nu}$,
establishing the non-renormalizability of GR.

\subsection{Stelle (1977)}

Stelle~\cite{Stelle1977} showed that the $R + R^2 + R_{\mu\nu}^2$
theory is power-counting renormalizable. The propagator scales as $1/k^4$
at large momenta, giving logarithmic one-loop divergences. However, the
theory contains a spin-2 ghost with unit negative residue.

%%======================================================================
\section{Prescription Uniqueness}
%%======================================================================

Anselmi, Briscese, Calcagni, and Modesto~\cite{ABCM2025}
(arXiv:2503.01841, 2025) analyzed four prescriptions for theories
with complex poles:
\begin{enumerate}
  \item \textbf{Textbook Wick rotation:} violates optical theorem
  \item \textbf{Lee-Wick-Nakanishi:} violates Lorentz invariance
  \item \textbf{Fakeon/AP:} violates analyticity (acceptable), preserves
        everything else
  \item \textbf{Direct Minkowski:} violates optical theorem AND power
        counting
\end{enumerate}

Only the fakeon prescription is physically viable for theories with
complex conjugate poles, including SCT.

%%======================================================================
\section{Nonlocal Gauge-Fixing}
%%======================================================================

\subsection{Calcagni--Modesto (2024)}

Calcagni and Modesto~\cite{CalcagniModesto2024}
(arXiv:2402.14785) derived the full gauge-fixed graviton propagator
in Barnes--Rivers basis for general nonlocal gravity, including:
\begin{itemize}
  \item Explicit treatment of gauge-fixing parameters (de~Donder,
        Julve--Tonin, Feynman gauges)
  \item Faddeev--Popov sector for nonlocal theories
  \item Convention dictionary between different references
\end{itemize}

\subsection{Tomboulis (1997) and Modesto (2012)}

Tomboulis~\cite{Tomboulis1997} (hep-th/9702146) pioneered nonlocal
gauge-fixing with $S_{\mathrm{gf}} = \tfrac{1}{2\alpha}\int F_\mu\,
G^{\mu\nu}(\Box/\Lambda^2)\, F_\nu$, where $G$ is entire.
Modesto~\cite{Modesto2012} (arXiv:1202.3151) developed the
super-renormalizable nonlocal gravity framework with exponential
entire functions.

%%======================================================================
\section{SCT-Specific References}
%%======================================================================

\begin{longtable}{@{}lll@{}}
\toprule
Document & Location & Content \\
\midrule
NT-4a linearized & theory/derivations/ & Propagator, Barnes--Rivers \\
PPN-1 derivation & theory/derivations/ & Newtonian potential \\
MR-2 derivation & theory/derivations/ & Ghost catalogue, conditions \\
OT verification & theory/consistency-checks/ & Optical theorem, $C_m = 283/120$ \\
KK resolution & theory/consistency-checks/ & Fakeon primary resolution \\
GP verification & theory/consistency-checks/ & Sign chain, width \\
\bottomrule
\end{longtable}

%%======================================================================
\section*{Acknowledgements}
%%======================================================================

workflow support.

%%======================================================================
\begin{thebibliography}{99}
%%======================================================================

\bibitem{ModestoCalcagni2021}
  L.~Modesto and G.~Calcagni,
  arXiv:2107.04558 [hep-th] (2021).

\bibitem{DonaGiaccari2015}
  P.~Dona et al.,
  arXiv:1506.04589 [hep-th] (2015).

\bibitem{AnselmiPiva2018a}
  D.~Anselmi and M.~Piva,
  arXiv:1803.07777 [hep-th] (2018).

\bibitem{ABCM2025}
  D.~Anselmi et al.,
  arXiv:2503.01841 [hep-th] (2025).

\bibitem{CalcagniModesto2024}
  G.~Calcagni and L.~Modesto,
  arXiv:2402.14785 [hep-th] (2024).

\bibitem{Stelle1977}
  K.~S.~Stelle,
  Phys.\ Rev.\ D \textbf{16}, 953 (1977).

\bibitem{tHooftVeltman1974}
  G.~'t~Hooft and M.~Veltman,
  Ann.\ Inst.\ H.\ Poincar\'e \textbf{20}, 69 (1974).

\bibitem{GoroffSagnotti1986}
  M.~H.~Goroff and A.~Sagnotti,
  Nucl.\ Phys.\ B \textbf{266}, 709 (1986).

\bibitem{Tomboulis1997}
  E.~T.~Tomboulis,
  hep-th/9702146 (1997).

\bibitem{Modesto2012}
  L.~Modesto,
  arXiv:1202.3151 [hep-th] (2012).

\end{thebibliography}

\end{document}
